\sect{त्रिस्पृशाख्यानम्}
\label{sec:padma-trisprisha}

\ganapatyadiDhyanam
\padmaPuranam

\dnsub{अथ षट्त्रिंशोऽध्यायः॥३६}

\uvacha{नारद उवाच}

\twolineshloka
{त्रिस्पृशाख्यं व्रतं ब्रूहि सर्वेश्वर विशेषतः}
{यच्छ्रुत्वा मुच्यते लोकः कर्मबन्धनतः क्षणात्}% ॥१॥

\uvacha{महादेव उवाच}

\twolineshloka
{सर्वपापौघशमनं महादुःखविनाशनम्}
{शृणु कृष्णावतारं त्वं त्रिस्पृशाख्यं महाव्रतम्}% ॥२॥

\twolineshloka
{कामदं सस्पृहाणां तु निस्पृहाणां तु मोक्षदम्}
{त्रिस्पृशाख्यं व्रतं विप्र शृणुष्व गदतो मम}% ॥३॥

\twolineshloka
{प्रत्यक्षमर्चितस्तेन कलिकाले च केशवः}
{त्रिस्पृशा कीर्तनं नित्यं यः करोति महामुने}% ॥४॥

\twolineshloka
{न पुरश्चरणे चीर्णे सर्वपापक्षयो भवेत्}
{त्रिस्पृशा नाममात्रेण क्षीयते नात्र संशयः}% ॥५॥

\twolineshloka
{नागमैर्न पुराणाद्यैर्न मखैस्तीर्थकोटिभिः}
{बहुभिर्व्रतसङ्घैश्च पूजितैस्त्रिदशैरपि}% ॥६॥

\twolineshloka
{मोक्षो भवति विप्रेन्द्र त्रिस्पृशा न कृता यदि}
{मोक्षार्थे देवदेवेन दृष्टा वै वैष्णवी तिथिः}% ॥७॥

\twolineshloka
{द्विजानां दुर्विदं साङ्ख्यं कलिकाले विशेषतः}
{अनिग्रहश्चेन्द्रियाणां स्थिरत्वं मनसो नहि}% ॥८॥

\twolineshloka
{विषयैर्विप्रयुक्तानां ध्यानधारणवर्जिनाम्}
{कामभोगप्रसक्तानां त्रिस्पृशा मोक्षदायिनी}% ॥९॥

\twolineshloka
{मह्यं चैव पुरा प्रोक्ता चतुर्वक्त्रस्य सागरे}
{क्षीरोदे प्रणतानां तु मत्स्यरूपेण चक्रिणा}% ॥१०॥

\twolineshloka
{त्रिस्पृशां ये करिष्यन्ति विषयैरपि संयुताः}
{तेषामपि मया दत्तो मोक्षः साङ्ख्यविवर्जिनाम्}% ॥११॥

\twolineshloka
{कामभोगप्रसक्तानां त्रिस्पृशा मोक्षदायिनी}
{बहुभिर्मुनिसङ्घैश्च कृतेयं च महामुने}% ॥१२॥

\twolineshloka
{कार्तिके शुक्लपक्षे तु त्रिस्पृशा जायते यदि}
{सोमेन सोमजेनापि पापकोटिविनाशिनी}% ॥१३॥

\twolineshloka
{यस्या उपोषणकृतो हत्यायुक्त महेशितुः}
{हस्ताद्ब्रह्मकपालं तु तत्क्षणात्पतितं भुवि}% ॥१४॥

\twolineshloka
{कलिकल्मषकोट्यौघैर्मुक्ता देवी त्रिमार्गगा}
{उपदेशान्माधवस्य त्रिस्पृशा समुपोषणात्}% ॥१५॥

\twolineshloka
{हत्याष्टौ बाहुवीर्यस्य पूर्वजाता महामुने}
{गता भृगूपदेशेन त्रिस्पृशा समुपोषणात्}% ॥१६॥

\twolineshloka
{शतायुधेन विप्रेन्द्र निहतो ब्राह्मणो वने}
{ब्रह्महत्याविनिर्मुक्तः त्रिस्पृशा समुपोषणात्}% ॥१७॥

\twolineshloka
{जीवोपदेशाच्छक्रस्य हत्या नमुचिसम्भवा}
{विनष्टा मुनिमुख्येन्द्र त्रिस्पृशा समुपोषणात्}% ॥१८॥

\twolineshloka
{ब्रह्महत्यादि पापानि त्रिस्पृशासमुपोषणात्}
{विलयं यान्ति विप्रेन्द्र पापेष्वन्येषु का कथा}% ॥१९॥

\twolineshloka
{न प्रयागे न काश्यां तु गोमत्यां कृष्णसन्निधौ}
{मोक्षो भवति विप्रेन्द्र त्रिस्पृशा यदि नो कृता}% ॥२०॥

\twolineshloka
{मरणाच्च प्रयागे तु गोमत्यां कृष्णसन्निधौ}
{स्नानमात्रेण गोमत्यां मुक्तिर्भवति शाश्वती}% ॥२१॥

\twolineshloka
{गृहेपि जायते मुक्तिस्त्रिस्पृशा समुपोषणात्}
{विषये वर्तमानस्य कामभोगान्वितस्य च}% ॥२२॥

\twolineshloka
{निवृत्तविषयस्यापि मुक्तिः साङ्ख्येन दुर्लभा}
{तस्मात्कुरुष्व विप्रेन्द्र त्रिस्पृशां मोक्षदायिनीम्}% ॥२३॥

\uvacha{नारद उवाच}

\twolineshloka
{कीदृशं तत्सुरश्रेष्ठ त्रिस्पृशाख्यं महाव्रतम्}
{मुक्तिदं यद् द्विजातीनां त्वया प्रोक्तं ममाधुना}% ॥२४॥

\uvacha{महादेव उवाच}

\twolineshloka
{जाह्नव्यै सा पुरा विप्र त्रिस्पृशा माधवेन तु}
{प्राची सरस्वती तीरे कथिता त्वनुकम्पया}% ॥२५॥

\uvacha{जाह्नव्युवाच}

\twolineshloka
{कलिकल्मष कोट्यौघैर्ब्रह्महत्यादिकैर्युताः}
{कलिकाले हृषीकेश स्नानं कुर्वन्ति मज्जले}% ॥२६॥

\twolineshloka
{तेषां पापशतैर्दोषैर्देहो मे कलुषी कृतः}
{कथं यास्यति मे देव पातकं गरुडध्वज}% ॥२७॥

\uvacha{प्राचीमाधव उवाच}

\twolineshloka
{कथयामि न सन्देहो मा पुत्रि रोदनं कुरु}
{श्यामोवटस्तु मे स्थानं प्राचीदेवी ममाग्रतः}% ॥२८॥

\twolineshloka
{वहते ब्रह्मतनया दृष्ट्वाग्रे च सुरेश्वरीम्}
{स्नानं कुरुष्व नित्यं त्वं त्वत्र पूता भविष्यसि}% ॥२९॥

\twolineshloka
{यत्र ब्रह्मसुता प्राची तत्राहं नात्र संशयः}
{तीर्थकोटिशतैर्युक्तः सुरैः सह वसाम्यहम्}% ॥३०॥

\twolineshloka
{पवित्रं मत्प्रियं स्थानं हत्याकोटिविनाशनम्}
{सन्तुष्टेन मया दत्तं यस्मात्प्राणाधिकासि मे}% ॥३१॥

\twolineshloka
{तीर्थकोटिसहस्राणि नित्यं तिष्ठन्ति जाह्नवि}
{प्राचीसरस्वती तोये सर्वदैव ममाज्ञया}% ॥३२॥

\twolineshloka
{ब्रह्मवधात्सुरापानात्गोधवाद्वृषलीवधात्}
{ब्रह्मस्वहरणादेव मातापित्रोस्त्वपूजनात्}% ॥३३॥

\twolineshloka
{चक्रियानाद्गुरुद्रोहादभक्ष्यस्य च भक्षणात्}
{सर्वपापस्य करणात्प्राची ब्रह्मसुता सुते}% ॥३४॥

\twolineshloka
{व्यपोहयति पापानि सकृत्स्नानेन मेऽग्रतः}
{कुरु स्नानं सरिच्छ्रेष्ठे विपापा त्वं भविष्यसि}% ॥३५॥

\uvacha{जाह्नव्युवाच}

\twolineshloka
{नाहं शक्नोमि देवेश आगन्तुं नित्यमेव हि}
{कथं नश्यन्ति पापानि कथयस्वेह माधव}% ॥३६॥

\uvacha{प्राचीमाधव उवाच}

\twolineshloka
{न शक्नोषि यदागन्तुं नित्यमेव हि जाह्नवि}
{तदान्यत्सम्प्रवक्ष्यामि यस्मान्मत्पादसम्भवा}% ॥३७॥

\twolineshloka
{सरस्वत्यधिका या च तीर्थकोटिशताधिका}
{मखकोट्यधिका वापि व्रतदानाधिका च या}% ॥३८॥

\twolineshloka
{जपहोमाधिका या च चतुर्वर्गफलप्रदा}
{साङ्ख्ययोगाधिका या च त्रिस्पृशा क्रियतां शुभा}% ॥३९॥

\twolineshloka
{यस्मिन्मासे समायाति सिता च यदि वासिता}
{कर्तव्या सा सरिच्छ्रेष्ठे कृते पापात्प्रमुच्यते}% ॥४०॥

\uvacha{जाह्नव्युवाच}

\twolineshloka
{कीदृशी त्रिस्पृशा देव ममाख्याहि सुमाधव}
{ईदृशो महिमा यस्यास्त्वया प्रोक्तो ममाधुना}% ॥४१॥

\twolineshloka
{दशम्येकादशीभद्रा दिनैकस्मिन्यदा भवेत्}
{त्रिस्पृशा सा भवेद्देव वान्यथा वद मे प्रभो}% ॥४२॥

\uvacha{कृष्ण उवाच}

\twolineshloka
{आसुरी त्रिस्पृशा देवी या त्वया परिकीर्तिता}
{वर्जनीया प्रयत्नेन वृत्तिहीनो यथा पतिः}% ॥४३॥

\twolineshloka
{असुराणां तु सा प्रोक्ता आयुर्बलविनाशिनी}
{वर्जनीया प्रयत्नेन यथा नारी रजस्वला}% ॥४४॥

\twolineshloka
{स्वजातिं च परित्यज्य या गताधमजातिषु}
{सैव त्याज्यं विशेषेण दशमीयुक्तं हि मद्दिनम्}% ॥४५॥

\twolineshloka
{यथा रजस्वलासङ्गाद्दूष्यन्ते ज्ञानवर्जिताः}
{तथैव दशमीयुक्तं मद्दिनं दूषितं नृणाम्}% ॥४६॥

\twolineshloka
{हत्यायुतशतं हन्ति त्रिस्पृशा समुपोषिता}
{एकादशी द्वादशी च रात्रिशेषे त्रयोदशी}% ॥४७॥

\twolineshloka
{त्रिस्पृशा सा तु विज्ञेया दशमीसहिता न हि}
{कृत्वापराधं मुच्येत प्रायश्चित्ते कृते नरः}% ॥४८॥

\twolineshloka
{दशमीवेधजं दोषं न क्षमामि सुरापगे}
{भुक्तं हालाहलं तेन विषस्य भक्षणं कृतम्}% ॥४९॥

\twolineshloka
{दशमीमिश्रितं येन कृतमेकादशीव्रतम्}
{इति मत्वा न कर्तव्यं मद्दिनं दशमीयुतम्}% ॥५०॥

\twolineshloka
{जन्मकोटिकृतं पुण्यं सन्तानं याति सङ्क्षयम्}
{पातयेत्स्वकुलं स्वर्गान्नयते रौरवादिकम्}% ॥५१॥

\twolineshloka
{स्वदेहं शोधयित्वा तु कर्तव्यो मम वासरः}
{वृद्धौ त्याज्या विना वेधाच्छ्रवणादिषु संयुता}% ॥५२॥

\twolineshloka
{जन्मपुण्यं क्षयं याति एकादश्युपवासिनाम्}
{संवृद्धौ तु विशेषेण सन्देहे समुपस्थिते}% ॥५३॥

\onelineshloka
{ममाज्ञया च कर्तव्या द्वादशीवल्लभा मम}% ॥५४॥

\uvacha{जाह्नव्युवाच}

\twolineshloka
{करिष्येऽहं जगन्नाथ त्रिस्पृशां वचनात्तव}
{सर्वपापविनिर्मुक्ता भविष्यामि तवाज्ञया}% ॥५५॥

\uvacha{श्रीकृष्ण उवाच}

\twolineshloka
{स्वस्थानं गच्छ भद्रं ते न भीः कार्या कदाचन}
{तव देवि सरिच्छ्रेष्ठे न पापं सङ्क्रमिष्यति}% ॥५६॥

\twolineshloka
{स्नात्वा सरस्वतीतोये येऽर्चयित्वा च माधवम्}
{प्रणमन्ति जगन्नाथं ते यान्ति परमां गतिम्}% ॥५७॥

\uvacha{जाह्नव्युवाच}

\twolineshloka
{विधानं ब्रूहि मे ब्रह्मन्सर्वस्वेन करोम्यहम्}
{प्रसादयामि देवेशं दामोदरमनामयम्}% ॥५८॥

\uvacha{प्राचीमाधव उवाच}

\twolineshloka
{शृणु देवि प्रवक्ष्यामि त्रिस्पृशाया विधानकम्}
{यं श्रुत्वापि सरिच्छ्रेष्ठे मुच्यते पातकैर्नरः}% ॥५९॥

\twolineshloka
{पलेन च पलार्धेन तदर्धेनापि वापगे}
{प्रतिमा मम सौवर्णा कार्या विभवसारतः}% ॥६०॥

\twolineshloka
{पात्रं ताम्रमयं कार्यं तिलैस्तु परिपूरितम्}
{सजलं तु घटं शुभ्रं पञ्चरत्नसमन्वितम्}% ॥६१॥

\twolineshloka
{वेष्टितं पुष्पमालाभिः कर्पूरागुरुवासितम्}
{न्यसेद्दामोदरं पश्चात्स्नापयित्वा विलिप्य च}% ॥६२॥

\threelineshloka
{परिधानं ततः कार्यं वस्त्रयुग्मेन चान्वितम्}
{मन्त्रैस्तु पूजनं कार्यं पुराणैः समुदीरितैः}
{पुष्पैः कालोद्भवैः शुभ्रैस्तुलसीपत्रैश्च कोमलैः}% ॥६३॥

\twolineshloka
{छत्रं तु विष्णवे दद्यात्पादुकाभ्यां सुसंयुतम्}
{निवेद्यानि मनोज्ञानि फलानि सुबहून्यपि}% ॥६४॥

\twolineshloka
{उपवीतं तु दातव्यं सोत्तरीयं नवं दृढम्}
{वैणवं दापयेद्दण्डं सुरूपं सोन्नतं दृढम्}% ॥६५॥

\twolineshloka
{दामोदराय वै पादौ जानुनी माधवाय च}
{गुह्यं कामप्रदायेति कटिं वामनमूर्तये}% ॥६६॥

\twolineshloka
{पद्मनाभाय नाभिं तु जठरं विश्वयोनये}
{हृदयं ज्ञानगम्याय कण्ठं वैकुण्ठगामिने}% ॥६७॥

\twolineshloka
{सहस्रबाहवे बाहू चक्षुषी योगरूपिणे}
{सम्पूज्य विधिवद्भक्त्या दद्यादर्घं विधानतः}% ॥६८॥

\twolineshloka
{शुभ्रेण नालिकेरेण शङ्खोपरि स्थितेन हि}
{सूत्रैरावेष्टितेनैव हस्तयोरुभयोरपि}% ॥६९॥

\twolineshloka
{स्मृतो हरसि पापानि यदि नित्यं जनार्दन}
{दुःस्वप्नं दुर्निमित्तानि मनसा दुर्विचिन्विता}% ॥७०॥

\twolineshloka
{नारकं तु भयं देव भयदुर्गतिसम्भवम्}
{यन्मम स्यान्महादेव ऐहिकं पारलौकिकम्}% ॥७१॥

\twolineshloka
{तेन देवेश मां रक्ष गृहाणार्घं नमोस्तु ते}
{कृपादृष्टिः सदैवास्तु दामोदर ममोपरि}% ॥७२॥

\twolineshloka
{धूपं दीपं च नैवेद्यं कुर्यान्नीराजनं ततः}
{शीर्षोपरि सरिच्छ्रेष्ठे भ्रामयेद्वारिजं हरेः}% ॥७३॥

\twolineshloka
{कृत्वा विधानमेतद्धि पूजयेत्स्वगुरुं ततः}
{दद्यात्सुवर्णं वस्त्राणि सोष्णीषं चैव कञ्चुकम्}% ॥७४॥

\twolineshloka
{उपानहोऽपि छत्रं च मुद्रिकां च कमण्डलुम्}
{भोजनं चैव ताम्बूलं सप्तधान्यं च दक्षिणाम्}% ॥७५॥

\twolineshloka
{गुरुं सम्पूज्य देवेशं कुर्याज्जागरणं हरेः}
{गीतनृत्यसमायुक्तं तथा शास्त्रसमन्वितम्}% ॥७६॥

\twolineshloka
{निशान्ते चैव देवाय दत्वा चार्घं विधानतः}
{स्नानादिकां क्रियां कृत्वा भुञ्जीयाद्वाडवैः सह}% ॥७७॥

\uvacha{शिव उवाच}

\twolineshloka
{द्विजैतत्त्रिस्पृशाख्यानमद्भुतं रोमहर्षणम्}
{श्रुत्वा तु लभते पुण्यं गङ्गां स्नानसमुद्भवम्}% ॥७८॥

\twolineshloka
{अश्वमेधसहस्राणि वाजपेयशतानि च}
{तत्फलं समवाप्नोति त्रिस्पृशा समुपोषणात्}% ॥७९॥

\twolineshloka
{पितृपक्षो मातृपक्षस्तथा चैवात्मपक्षकः}
{तैः सर्वैः सह सम्मुक्तो विष्णुलोके महीयते}% ॥८०॥

\twolineshloka
{तीर्थकोटिषु यत्पुण्यं क्षेत्रकोटिषु यत्फलम्}
{तत्फलं समवाप्नोति त्रिस्पृशा समुपोषणात्}% ॥८१॥

\twolineshloka
{ब्राह्मणा येऽपि कुर्वन्ति क्षत्रियाः कृष्णमानसाः}
{वैश्या वा शूद्रजन्मानो ये तथा चान्यजातयः}% ॥८२॥

\twolineshloka
{ते सर्वे मुक्तिमायान्ति भुवं त्यक्त्वा द्विजोत्तम}
{मन्त्राणां मन्त्रराजोऽथ यथा स्याद्द्वादशाक्षरः}% ॥८३॥

\twolineshloka
{व्रतानां च यथा चैषा येन वै त्रिस्पृशा कृता}
{ब्रह्मणा च कृता पूर्वं पश्चाद्राजर्षिभिः कृता}% ॥८४॥

\twolineshloka
{अन्येषां का कथा वत्स त्रिस्पृशा मुक्तिदायिनी}
{अनेन विधिना ब्रह्मंस्त्रिस्पृशासम्भवं व्रतम्}% ॥८५॥

\twolineshloka
{यः करोति नरो भक्त्या शृणु वक्ष्यामि तत्फलम्}
{गङ्गावगाहने ब्रह्मन्वाराणस्यां तु यत्फलम्}% ॥८६॥

\twolineshloka
{मन्वन्तरसहस्रैस्तु त्रिस्पृशाकारको हि तत्}
{प्राची च यमुना स्नाने वर्षैर्यत्कोटिभिः फलम्}% ॥८७॥

\twolineshloka
{तत्फलं समवाप्नोति त्रिस्पृशाव्रतकृन्नरः}
{तत्फलं तु कुरुक्षेत्रे सूर्यग्रहणकोटिभिः}% ॥८८॥

\twolineshloka
{हेमभारशतैर्दानैस्त्रिस्पृशाकरणेन तत्}
{पापकोटिसहस्राणि हत्याकोटिशतानि च}% ॥८९॥

\twolineshloka
{एकेनैवोपवासेन क्रियते भस्मसाद्द्रुतम्}
{त्रिस्पृशाया व्रतं यत्तु अगतीनां गतिप्रदम्}% ॥९०॥

\twolineshloka
{गतिमिच्छन्ति विप्रर्षे महत्पापशतानि च}
{स्वयं कृष्णेन कथितं पाराशर्यस्य चाग्रतः}% ॥९१॥

\twolineshloka
{प्रकाशयति यश्चेदं लिखित्वा वैष्णवं द्विजे}
{पापोघैर्ग्रथितस्यापि तस्य मुक्तिर्भविष्यति}% ॥९२॥

\twolineshloka
{पुण्यैरवाप्यते विद्वन्मन्वन्तरशतैरपि}
{त्रिस्पृशा दुर्लभा लोके प्राप्यते नैव मानवैः}% ॥९३॥

\twolineshloka
{कलौ ये त्रिस्पृशां लब्ध्वा न कुर्वन्ति नराधमाः}
{तेषां जन्मफलं चैव जीवितं विफलं भवेत्}% ॥९४॥

\twolineshloka
{प्रेतत्वं तैः समुत्तीर्णं विना श्राद्धैर्विना सुतैः}
{कृता यैस्त्रिस्पृशा विद्वन्सकृत्प्राप्य कलौ युगे}% ॥९५॥

॥इति श्रीपाद्मे महापुराणे पञ्चपञ्चाशत्साहस्र्यां संहितायामुत्तरखण्डे उमापतिनारदसंवादे त्रिस्पृशाख्यानं नाम षट्त्रिंशोऽध्यायः॥३६॥

\genMangalaShloka

\hyperref[sec:ekadashi_mahatmyam_padma_puranam]{\closesub}
\clearpage

